\documentclass[a4paper,10pt]{article}
\usepackage{multicol}
\usepackage{epsfig,graphicx,subfigure,amsthm,amsmath}
\usepackage{color,xcolor}
\usepackage{fancyhdr}
\usepackage[top=3cm,bottom=5cm,right=2.5cm,left=2.5cm]{geometry}       
\usepackage{xepersian}
\pagestyle{fancy}
\chead{\includegraphics[width=455pt]{logo_imc.jpg}}
\lhead{}
\rhead{}
\renewcommand{\headrulewidth}{.2cm}
%\fancyheadoffset[RE,LO]{0\textwidth}
\renewcommand{\headrule}{\color{orange}\hrule width\headwidth  height\headrulewidth \vskip-\headrulewidth}
\setlength{\headsep}{1cm}
\pagenumbering{arabic} \setcounter{page}{1}
\label{sec.sub.gdn}

\settextfont[Scale=1.2]{XB_Niloofar.ttf}
\setlatintextfont[Scale=1]{Times New Roman}
\renewcommand{\baselinestretch}{1.5}


\title{\textbf{ چکیده مبسوط مقالات }}
\date{}
\author{\small محمدرضا صافی ، عضو هيأت علمی گروه ‌ریاضی، دانشگاه سمنان، msafi@profs.semnan.ac.ir\\
\small نفر دوم، ....\\
\small نفر سوم، .....\\
\small  (ذكرسمت مؤلفين شامل عضو هيأت علمی، دانشجوی دكتری، دانشجوی كارشناسی ارشد، \\
\small فارغ التحصيل دكتری، فارغ التحصيل كارشناسی ارشد و غيره در اين قسمت ضروری است)}


\begin{document}

\maketitle
\thispagestyle{fancy}

\rule{\textwidth}{0.2mm}

\noindent
\textbf{چکیده:}
 در این نوشته نحوه نگارش چکیده مبسوط مقالات ارائه شده است. این چکیده بایستی حداکثر شامل 200 کلمه باشد و تا حد امکان از فرمولها و معادلات در آن صرفنظر شود. نویسندگان بایستی در یک یا دو پاراگراف نتایج خود را بیان نمایند. \\
\textbf{ کلمات کلیدی:}
 نویسندگان حداکثر 5 کلمه از موضوعات مهم مقاله خود را انتخاب و در این بخش قرار دهند. 
\rule{\textwidth}{0.2mm}


\setlength{\columnsep}{1cm}
\begin{multicols}{2}

\section*{مقدمه}

در این نوشته نحوه نگارش چکیده مبسوط مقالات چهل و پنجمین کنفرانس ریاضی ایران به اختصار بیان می‌شود تا علاقمندان به شرکت در کنفرانس به راحتی مقالات خود را به فرمت موردنظر کنفرانس درآورند. 

\section*{بخش یک}

در این بخش فونت مورد استفاده برای آماده‌سازی مقالات فارسی کنفرانس را بیان می‌کنیم. در متن مقاله از فونت
\lr{ XB Niloofar}
  استفاده شده است. 

\section*{بخش دو}

شرکت‌کنندگان بایستی مقالات خود را با استفاده از الگوی ارائه‌شده در وبسایت کنفرانس آماده نمایند و حداکثر تا تاریخ 1393/2/31 به دبیرخانه کنفرانس ارسال کنند. 

\section*{بخش سه}

نتایج مقالات پس از بررسی در کمیته علمی کنفرانس در وبسایت اطلاع‌رسانی خواهد شد. همچنین از علاقمندان به برگزاری کارگاهها درخواست می‌شود موضوعات پیشنهادی خود را در اسرع وقت به دبیرخانه کنفرانس ارسال نمایند تا درخصوص برگزاری آنها تصمیم‌گیری شود. 

\section*{نتایج}

کلیه شرکت‌کنندگان در همایش می‌بایستی حداکثر تا تاریخ 1393/3/31 ثبت‌نام بعمل آورند. 



\begin{thebibliography}{99}

\bibitem{}

ر. حیدری، رفتارگرایی و طرح درس در بوته نقد، مجله رشد آموزش ریاضی (1383)، شماره 78، 23-32

\begin{LTRitems}

\resetlatinfont

\bibitem{abramsky2}
S. Abramsky, A. Jung, {\em Domain theory}, in: S. Abramsky, D.M. Gabbay, T.S.E. Maibaum (Eds.), Handbook of
Logic in Computer Science, Vol. 3, Clarendon Press, Oxford, 1994, pp. 1–68.

\bibitem{aliprantis}
C.D. Aliprantis and O. Burkinshaw,   {\em Principles of Real Analysis}.  Academic Press.  1998,  xii+415 pp.

\bibitem{Alvarez1}
M. Alvarez-Manilla, {\em Measure theoretic results for continuous valuations on partially ordered
spaces}, Dissertation, Imperial College, London, 2000.


\end{LTRitems}

\end{thebibliography}


\end{multicols}


\end{document}